\pdfoutput=1
\ifdefined\pdftrailerid\pdftrailerid{}\fi
\documentclass[11pt]{article}
\usepackage[T1]{fontenc}
\usepackage{lmodern,microtype}
\usepackage[margin=.95in]{geometry}
\usepackage{amsmath,amssymb,amsthm,mathtools}
\usepackage[colorlinks=true,allcolors=blue!55!black]{hyperref}
\usepackage{xcolor}
\newtheorem{theorem}{Theorem}[section]
\newtheorem{proposition}[theorem]{Proposition}
\newtheorem{lemma}[theorem]{Lemma}
\newtheorem{corollary}[theorem]{Corollary}
\theoremstyle{remark}
\newtheorem{remark}[theorem]{Remark}
\newcommand{\dd}{\mathrm d}
\newcommand{\R}{\mathbb R}
\newcommand{\scrL}{\mathcal L}
\newcommand{\eps}{\varepsilon}
\setlength{\emergencystretch}{1em}
\allowdisplaybreaks[1]
\title{Vacuum stress spectroscopy with double pulses\\ and two-scale thresholds}
\author{}
\date{}
\hypersetup{pdftitle={Vacuum stress spectroscopy with double pulses and two-scale thresholds},pdfauthor={}}
\begin{document}
\maketitle
\begin{abstract}
We reconstruct the connected Wightman two-point function of a specified
conserved stress tensor in four-dimensional Minkowski space from a
restricted family of vacuum variance measurements. A single compact
spatial profile and one compact time pulse are fixed; the time pulse is
used twice with a variable delay. The energy-density scan determines
a positive mass-squared measure including its null atom, and a trace or
spatial transverse-traceless scan separates the timelike spin channels.
The exact inverse is not uniformly stable: free Klein--Gordon stresses
with the same linear curvature improvement have arbitrarily close
fixed-window scans but different null residues. We derive sharp
asymptotics when spatial size and time duration both vary, including
the slowly varying transition at threshold index \(1/2\). A converse
recovers the threshold index and normalization from the regular
variation of the reference scan and convergence of the delay ratios.
\end{abstract}

\section{Stress data and measurement model}

A positive-energy stress Wightman function has a spectral
decomposition fixed by covariance and conservation up to two
positive timelike measures \cite{CFL,Karateev}. If access is limited
to compact stress windows, it is not evident whether the measure is
identifiable or whether a measured infrared power is a property of the
detector rather than the threshold. We keep the stress prescription,
detector calibration, and state fixed. All data below are vacuum norms
of smeared stress vectors. Neither a time-ordered correlator nor an
independent measurement of a Fourier phase is needed.

Use signature \((-+++)\), \(p=(\omega,\mathbf p)\),
\(q=|\mathbf p|\), and \(s=-p^2\).
Let \(T_{\mu\nu}\) be a centered, Hermitian, symmetric, conserved
operator-valued tempered distribution in a positive Hilbert space
with a Poincar\'e invariant vacuum vector and spectrum in the future
cone. Throughout the main results the centered vacuum columns have
no component at zero \emph{four}-momentum; this is automatic in a
unique invariant vacuum after centering. We use the Fourier inverse
measure \(\dd^d k/(2\pi)^d\). In the timelike region let
\[
 D_{\mu\nu}=p_\mu p_\nu+s\eta_{\mu\nu},\qquad
 P^{(0)}_{\mu\nu,\rho\sigma}=\tfrac13D_{\mu\nu}D_{\rho\sigma},
\]
\[
 P^{(2)}_{\mu\nu,\rho\sigma}=
 \tfrac12(D_{\mu\rho}D_{\nu\sigma}+D_{\mu\sigma}D_{\nu\rho})
 -\tfrac13D_{\mu\nu}D_{\rho\sigma}.
\]
The nonnegative measures \(\rho_0,\rho_2\) are in the normalization
where the rest-frame factors \(s^2\) are included in these tensors.
Indeed in the rest frame conservation removes time components,
and the symmetric spatial tensor decomposes into its trace and
traceless irreducible parts. Positivity of the Wightman matrix
measure fixes the sign of each spectral weight.

The punctured future null cone contributes only a nonnegative
coefficient \(a\) times \(p_\mu p_\nu p_\rho p_\sigma\) with the
invariant shell measure \cite{P45}. Recall the reason: on that
homogeneous orbit the positive matrix range carries a unitary
action of the little group. Null translations, unipotent in the
finite tensor representation, must act as identity on the range.
Conservation leaves the symmetric square of the span of \(p\) and
two transverse vectors \(e_i\). Under the independent shifts
\(e_i\mapsto e_i+c_i p\), invariance first eliminates the transverse
quadratic terms and then the mixed terms, leaving \(p\otimes p\).
Thus the energy contraction is
\begin{equation}
 \widehat W_{00,00}(\omega,\mathbf p)
 =2\pi\theta(\omega)q^4\int_{[0,\infty)}
        \delta(\omega^2-q^2-s)\,\nu(\dd s),\qquad
 \nu=a\delta_0+\tfrac13(2\rho_2+\rho_0)|_{s>0}.
 \label{eq:wightman}
\end{equation}
The origin atom excluded above is different from \(a\). All these
mass-squared measures are locally finite and have a finite
polynomially weighted tail. For \(\nu\), this follows by testing
\eqref{eq:wightman} on a compact spatial annulus with \(q>0\) and
frequencies covering the indicated mass interval. The shell weight
is bounded below on a smaller such annulus. At large frequency
take a sufficiently decaying polynomial test and use temperateness
of the positive matrix measure. Positivity of the two contributions
then gives the same bounds on \(\rho_j\) on \(s>0\).

Fix real \(0\ne w\in C_c^\infty(\R^3)\), and real
\(g\in C_c^\infty(\R)\) with \(\int g=1\). For a time gate \(f\)
write \(V_\nu(f;w)\) for the squared norm of
\(\int f(t)w(\mathbf x)T_{00}(t,\mathbf x)\dd t\,\dd^3x\)
applied to the vacuum. In the first experiment neither \(w\) nor the
shape \(g\) is changed. At delay \(\theta\) we use
\[
            g^{[\theta]}(t)=\tfrac12[g(t)+g(t-\theta)].
\]
These are actual normalized compact gates. Below we also dilate
their physical duration and spatial aperture.

\section{Exact inversion with one fixed spatial profile}

For a nonnegative integer \(j\) used in a given detector, define
the known positive spatial measure
\[
 \eta_{w,j}=(q^2)_*\bigl(q^j |\widehat w(\mathbf p)|^2
                              \dd^3\mathbf p/(2\pi)^3\bigr).
\]
Only \(j=0,4\) will be needed. For \(r>0\) its density is
\begin{equation}
 h_{w,j}(r)=\frac{r^{(j+1)/2}}{2(2\pi)^3}
          \int_{S^2}|\widehat w(\sqrt{r}\,\vartheta)|^2\,\dd\vartheta .
 \label{eq:spatialdensity}
\end{equation}
This measure is nonzero, absolutely continuous, and rapidly decaying
at infinity. In particular it has no atom at the origin.
Directly integrating the positive shell in \eqref{eq:wightman}
gives the energy frequency measure \(\xi_\nu\):
\begin{equation}
 V_\nu(f;w)=\int_0^\infty |\widehat f(\omega)|^2\,\xi_\nu(\dd\omega),
 \qquad
 (\omega^2)_*\xi_\nu(\dd u)=
       \frac{1}{2\sqrt u}(\eta_{w,4}*\nu)(\dd u).
 \label{eq:detector}
\end{equation}
The shell factor \(1/(2\sqrt u)\) appears before the pushforward;
there is no further Jacobian when \((\omega^2)_*\) acts on measures.

\begin{theorem}[Fixed-window inverse]
\label{thm:inverse}
For the stated \(g,w\), the exact values \(V_\nu(g^{[\theta]};w)\)
at a countable dense set of nonnegative delays together with
\(V_\nu(g;w)\) determine \(\nu\). This includes the null residue,
positive-mass atoms, gaps, and singular continuous components.
The same inversion applies to \eqref{eq:detector} when \(4\)
is replaced by \(0\) and \(\nu\) by any positive locally finite
mass measure with a polynomially weighted finite tail.
\end{theorem}

\begin{proof}
The temporal Fourier modulus gives
\begin{equation}
 2V_\nu(g^{[\theta]};w)-V_\nu(g;w)=
        \int_0^\infty \cos(\theta\omega)\,\beta_\nu(\dd\omega),
 \qquad \beta_\nu=|\widehat g|^2\xi_\nu .
 \label{eq:cosine}
\end{equation}
The right side is continuous in \(\theta\) and is the Fourier
transform of the even symmetrization of a finite positive measure.
Thus the data recover \(\beta_\nu\). Its finiteness follows from
the spectral tail and the rapid temporal decay.

The zeros of \(\widehat g\) on the real axis are isolated, since
\(\widehat g\) is entire and \(\widehat g(0)=1\). They lose no
measure information. In fact \(\eta_{w,j}*\nu\) is locally
absolutely continuous by \eqref{eq:spatialdensity} and Tonelli,
even for an atomic \(\nu\). Near zero the convolution density is
\(O(u^{(j+1)/2})\), using local finiteness of \(\nu\).
Multiplication by \(u^{-1/2}\) is therefore locally integrable;
the resulting frequency measure remains absolutely continuous
and has no origin atom. Divide by \(|\widehat g|^2\) off its zeros,
push to \(u\), and multiply by \(2\sqrt u\). This recovers the
convolution of measures, denoted \(\Upsilon_\nu\).

For every \(p>0\) one can now form
\begin{equation}
       \scrL\nu(p)=
        \frac{\scrL\Upsilon_\nu(p)}{\scrL\eta_{w,j}(p)}.
 \label{eq:laplace}
\end{equation}
Both transforms in the quotient are finite; the denominator is
strictly positive. Exponential tilting reduces uniqueness of
\(\scrL\nu\) to finite-measure Fourier uniqueness: two Laplace
transforms agreeing on \(p>0\) agree in the open right half-plane
by analyticity, and the vertical line through any positive
\(p\) is the Fourier transform of the tilted measures.
\end{proof}

For example \(a=\lim_{p\to\infty}\scrL\nu(p)\).
For a nonzero remaining measure its lower support is
\(-\lim_{p\to\infty}p^{-1}\log(\scrL\nu(p)-a)\).
The latter fact follows by exponential tilting and bounding
above and below using any neighborhood of the lower support.
The unbounded delay family and exact readings are substantive.
Section~\ref{sec:noise} gives discontinuity for the inverse
question of the presence of the null atom.

\begin{remark}[Separating a possible zero-momentum component]
If a centered vacuum column can have an atom at \(p=0\), its
positive Lorentz-invariant tensor is \(b\,\eta_{\mu\nu}
\eta_{\rho\sigma}\) with \(b\ge0\). To see this, the finite
positive range carries a unitary Lorentz action. Each boost in
the symmetric tensor representation is diagonalizable with
positive real eigenvalues, hence acts trivially on this range.
Boost commutators also generate the rotations, leaving the
invariant metric line. For a nonzero spatial mean the energy
gate records this atom as \(b|\widehat w(0)|^2\delta_0\)
in \(\beta\) of \eqref{eq:cosine}; the null shell instead
produces an absolutely continuous frequency distribution.
Subtract the measured frequency atom before using
\eqref{eq:laplace}. A zero-mean spatial profile cannot see
this origin component. For the trace gate the corresponding
origin weight is sixteen times the energy weight.
\end{remark}

\section{Resolving the timelike spin channels}

The energy family determines \(2\rho_2+\rho_0\) but does not
separate the two measures by itself. A second fixed contraction
of the physical stress suffices.

\begin{proposition}[Two compact detector families suffice]
\label{thm:spin}
Under the stated hypotheses the \(T_{00}\) double-pulse scan,
together with either the trace double-pulse scan at the same
\(w\) or a scan with one fixed nonzero compact spatial
transverse-traceless tensor, determines both timelike measures
and \(a\). These data fix the entire specified connected stress
two-point function.
\end{proposition}

\begin{proof}
For the trace \(\Theta=\eta^{\mu\nu}T_{\mu\nu}\),
\(\eta^{\mu\nu}D_{\mu\nu}=3s\) and the spin-two tensor is
traceless. Its full double contraction with \(P^{(0)}\)
is \(3s^2\), and the null-pole tensor has zero trace.
Thus the trace scan has formula \eqref{eq:detector} with
spatial measure \(\eta_{w,0}\) and mass measure
\(3s^2\rho_0(\dd s)\). Reconstruct the latter using
Theorem~\ref{thm:inverse}, divide on \(s>0\), and use
\eqref{eq:wightman} to obtain \(\rho_2\).

For the other option take a real symmetric compact spatial tensor
\(f_{ij}\) such that \(\partial_i f_{ij}=0\),
\(\delta^{ij}f_{ij}=0\), with no time components in the test.
These profiles exist by applying the local polynomial tensor
operator with symbol
\[
 \mathcal P_{ij,ab}(\mathbf k)=
  \tfrac12(P_{ia}P_{jb}+P_{ib}P_{ja})-\tfrac12P_{ij}P_{ab},
 \qquad P_{ij}=|\mathbf k|^2\delta_{ij}-k_i k_j
\]
to a compact smooth template whose transverse-traceless part
is not identically zero. The identities
\(\mathbf k^iP_{ij}=0\), \(P^2=|\mathbf k|^2P\),
and \(\operatorname{tr}P=2|\mathbf k|^2\) give the constraints
and preserve compact support because the operator is differential.
Then \(D:\widehat f=0\) and direct contraction gives
\(\overline{\widehat f}:P^{(2)}:\widehat f
=s^2|\widehat f|^2\). Its spatial blur is formed from
\(|\widehat f|^2\), with the tensor norm summed over the indices,
and its mass measure is \(s^2\rho_2(\dd s)\).
The same inversion and division recover \(\rho_2\); the energy
record supplies \(\rho_0\) and the null residue.
\end{proof}

These are restricted-access measurements rather than a new
spectral decomposition. The distinction between the timelike
spin-zero measure and the null helicity-zero atom matters:
the trace is blind to the latter despite the name of the
timelike sector. For example the mixed energy--trace contraction
has weight \(q^2s\,\rho_0\), with no null response.

\begin{proposition}[An obstruction to a universal scalar-only null probe]
\label{thm:nullblind}
Let \(F_{\mu\nu}(p)\) be a detector tensor continuous to a fixed
nonzero future null \(p_0\), as for a compact smooth test. If
its timelike spin-two norm vanishes on the nearby timelike
orbit directions tending to \(p_0\) with fixed spatial momentum,
then \(F_{\mu\nu}(p_0)p_0^\mu p_0^\nu=0\).
Thus exact blindness of one such probe to arbitrary timelike
spin-two spectral data implies blindness to the null atom there.
\end{proposition}

\begin{proof}
Write \(p_0=(q,qn)\), \(q>0\), and
\(p_s=(\sqrt{q^2+s},qn)\). Positivity of the timelike projector
forces the transverse-to-\(p_s\) three-dimensional block of \(F\)
to be proportional to the identity. Evaluate it on a fixed
unit purely spatial \(e_1\perp n\) and on the spacelike unit
vector \(e_L=(q,\sqrt{q^2+s}\,n)/\sqrt s\) perpendicular to \(p_s\).
Their evaluations agree. Multiply by \(s\) and pass to the
limit. The right side from \(e_1\) tends to zero while
\(\sqrt s\,e_L\to p_0\), proving the claim.
\end{proof}
This does not preclude separating the null atom after combining
different scans. Nor does it preclude choosing a probe adapted to
a known spectrum with no timelike masses accumulating at zero.

\section{The fixed-window inverse has no uniform noise margin}
\label{sec:noise}

\begin{theorem}[Physical instability at a null threshold]
\label{thm:noise}
For any fixed \(w,g\) as above, even the two full delay scans
of the energy and the trace have no uniformly positive
noise margin for the presence of a nonzero null atom over
this stress class. There is a family of free Klein--Gordon
stresses for which both scans converge uniformly in all delays
as \(m\downarrow0\), whereas the null residue is positive at
\(m=0\) and zero at every \(m>0\).
\end{theorem}

\begin{proof}
First compare the norm of the gated measures that appear in
\eqref{eq:cosine}. In the spectral cone, replacing
\(a\delta_0\) by \(\rho_0=3a\delta_\eps\), \(\eps>0\),
changes the energy weight to \(a\delta_\eps\)
and the trace weight to \(9a\eps^2\delta_\eps\).
For a mass atom at \(\eps\), the direct squared-frequency
density of the gated measure, per unit mass weight, is
\[
 b_g(u)h_{w,j}(u-\eps),\qquad
       b_g(u)=\frac{|\widehat g(\sqrt u)|^2}{2\sqrt u} .
\]
The spatial density is taken to vanish at negative arguments.
These densities converge in \(L^1(\dd u)\) to the expression
at \(\eps=0\). Indeed \(h_{w,j}(u-\eps)=O(u^{(j+1)/2})\)
uniformly near zero; away from zero use continuity and rapid
temporal decay. Hence the displaced energy atom gives uniform
convergence of its whole scan: the double-pulse multiplier
\((1+\cos(\theta\sqrt u))/2\) has absolute value at most one,
independently of \(\theta\). The trace contribution vanishes
in the same norm with its additional \(\eps^2\) factor. A
positive mixture of masses in a vanishing interval above zero
gives the same instability, with no isolated positive-mass
atom required at the measure level.

This behavior occurs for actual local conserved stresses.
Let \(\varphi_m\) be a real free scalar in four dimensions,
and prescribe the same linear curvature improvement at each
mass:
\[
 T_{\mu\nu}^{(m)}=:T_{\mu\nu,\mathrm{can}}^{(m)}:
  +\ell(\eta_{\mu\nu}\Box-\partial_\mu\partial_\nu)\varphi_m,
 \qquad \ell\ne0 .
\]
The improvement is a background-curvature coupling and changes
the translation charge only by a spatial boundary term in the
usual test domain. The mixed Wightman term between the canonical
quadratic stress and this linear one vanishes by free-field
parity. On a one-particle shell the improvement contributes
\(\ell D_{\mu\nu}\) in canonical scalar normalization.
Consequently the energy mass measure includes precisely
\(\ell^2\delta_{m^2}\) from the one-particle term, and its
trace measure includes \(9\ell^2m^4\delta_{m^2}\).
At \(m=0\) this gives a null-shell residue \(a=\ell^2\).

The two-particle continuum from the canonical stress converges
in the same gated total-variation norm. For \(s\ge4m^2\), in the
pair rest frame the two momentum magnitudes are
\(\sqrt{s-4m^2}/2\) and their energies are \(\sqrt s/2\).
The quadratic stress form factor is a fixed polynomial of
degree two in the momenta and \(m\), bounded by \(Cs\);
after division by \(s\) it converges for every fixed \(s>0\).
The invariant two-particle phase space gives a fixed constant times
\(\sqrt{1-4m^2/s}\,\mathbf1_{s\ge4m^2}\) after angular integration.
Thus in the \(s^2\)-factored convention of
\eqref{eq:wightman} the positive canonical spin densities
are uniformly bounded, and converge pointwise for \(s>0\).
The trace density is bounded by \(Cs^2\). The weighted output
total variation of either density difference is bounded, by
Tonelli, by the integral of its absolute value against
\(\int_0^\infty b_g(r+s)\eta_{w,j}(\dd r)\).
This kernel is locally bounded, since
\(b_g(r+s)\le C/\sqrt{r}\) there, and is rapidly decreasing
at large \(s\). Dominated convergence proves the claimed
uniformity of the full records.
\end{proof}

For these physical examples one can also bound the approach
to the limiting record. The energy spatial blur density
extended by zero has bounded first derivative near zero
and at infinity: the explicit formula gives
\(h_{w,4}'(r)=O(r^{3/2})\) as \(r\downarrow0\).
The one-particle energy change is thus
\(O(m^2)\) in gated total variation, using the fundamental
theorem of calculus under \(b_g(u)\dd u\). The rest-frame
canonical contractions, divided by \(s^2\), are smooth
functions of \(m^2/s\) on \(0\le m^2/s\le1/8\): the
polynomial form factors contain spatial quadratic
momenta, and the phase-space square root is smooth on
this interval. On the remaining allowed region they
are uniformly bounded. Their differences from the
massless densities are consequently bounded by
\(C\min(1,m^2/s)\). Pairing with the same gated kernel gives
uniform errors \(O(m^2(1+|\log m|))\) in the energy scan
and \(O(m^2)\) in the trace scan as \(m\downarrow0\); the latter
uses the additional trace weight \(s^2\). The bounds are
only upper bounds and do not supply an optimal resolution.

The theorem excludes a uniformly correct null-atom decision
at any fixed positive worst-case error on the fixed-window
records. It does not say that a specified finite spectral family
is always hard to distinguish, or that scaling up the spatial
window cannot separate the cases. In particular it does not
conflict with exact uniqueness in Theorem~\ref{thm:inverse}.

\section{Two different threshold experiments}
\label{sec:thresholds}

We distinguish the threshold law of the recovered additive
convolution at a fixed window from the multiplicative scaling law
of the raw variance. These use different observations.

\subsection{The additive spatial blur at an arbitrary mass threshold}

Write \(\Upsilon=\eta_{w,j}*\mu\), where \(\mu\) is the mass measure
for one detector, with \(j=4\) for the energy and \(j=0\) for the
trace or compact spatial tensor. In the latter case use the
tensor norm in the spatial profile. Let \(s_0=\inf\operatorname{supp}
\mu<\infty\) for nonzero \(\mu\). If \(P_m\) is the first nonzero
homogeneous Taylor polynomial at zero of the spatial Fourier profile,
\eqref{eq:spatialdensity} gives
\[
 h_{w,j}(r)\sim c_{w,j} r^{\gamma-1},\qquad
 \gamma=(j+3)/2+m,\qquad
 c_{w,j}=\frac{1}{2(2\pi)^3}\int_{S^2}|P_m(\vartheta)|^2\dd\vartheta>0 .
\]
The finite minimal order \(m\) exists by analyticity.
Positivity therefore gives
\(\inf\operatorname{supp}\Upsilon=s_0\);
it is also the square of the lower frequency support of
the gated output, since isolated zeros of \(\widehat g\) carry
no output mass.

\begin{proposition}[The exact threshold shift of the additive blur]
\label{thm:additive}
For \(\alpha\ge0\), \(C>0\), and \(L\) positive and slowly
varying at infinity, the following two asymptotics are equivalent:
\begin{align*}
 \mu([s_0,s_0+x])&\sim Cx^\alpha L(1/x),\\
 \Upsilon([s_0,s_0+x])&\sim
 \frac{C c_{w,j}\Gamma(\gamma)\Gamma(\alpha+1)}
      {\Gamma(\alpha+\gamma+1)}
       x^{\alpha+\gamma}L(1/x).
\end{align*}
This includes an atom at \(s_0\).
\end{proposition}

\begin{proof}
Translate to \(s_0=0\). The Laplace--Stieltjes Karamata
equivalence \cite{Karamata,BGT} changes a positive cumulative
asymptotic \(Cx^\alpha L(1/x)\) into
\(C\Gamma(\alpha+1)p^{-\alpha}L(p)\), and conversely.
For clarity, the forward implication follows by positive
rescaling and tail bounds. For the converse rescale the
measure under \(x\mapsto px\) with factor \(p^\alpha/L(p)\).
Its Laplace transforms converge at every positive test
parameter, giving compact mass control and exponentially
weighted tails from half that parameter. Uniqueness of the
positive limiting Laplace transform gives the power-law
measure (or the origin atom for \(\alpha=0\)), so evaluation
at \([0,1]\) gives the cumulative constant.
Also
\(\scrL\eta_{w,j}(p)\sim c_{w,j}\Gamma(\gamma)p^{-\gamma}\)
by its explicit density. Applying the equivalence to
\(\scrL\Upsilon=\scrL\eta_{w,j}\scrL\mu\) proves the proposition
in both directions.
\end{proof}

\subsection{The balanced dilation law}

For the energy measurement now let
\(w_R(\mathbf x)=w(\mathbf x/R)\) and
\(g_\tau(t)=\tau^{-1}g(t/\tau)\).
Put \(b=\tau/R\). For \(j\ge0\) write
\(W_j=\int |\mathbf k|^j|\widehat w(\mathbf k)|^2\dd^3\mathbf k/(2\pi)^3\).
A direct change of variables in the exact shell integral gives
\begin{equation}
 V(R,\tau)=\int F_b(R^2s)\nu(\dd s),\qquad
 F_b(y)=\frac12\int\frac{\dd^3\mathbf k}{(2\pi)^3}
       \frac{|\mathbf k|^4|\widehat w(\mathbf k)|^2}{\sqrt{|\mathbf k|^2+y}}
       \left|\widehat g\bigl(b\sqrt{|\mathbf k|^2+y}\bigr)\right|^2 .
 \label{eq:scale}
\end{equation}
The factor \(R^3\) from the two spatial profiles and volume,
the \(R^{-4}\) from \(q^4\), and the \(R\) from the shell
denominator cancel. For each \(b>0\), \(F_b\) is continuous
and positive on \([0,\infty)\), with rapid decay at infinity.
The rapid bounds and local continuity are uniform on compact
subintervals of \(b>0\). Positivity holds even for a time
gate with notches: the frequency zero set is discrete and
the spatial integral is not supported on it.

For the following results let \(U(s)=\nu([0,s])\).
One can instead study only the continuum by removing the exact
atom term \(a F_b(0)\) from all the variances and using
\(U(s)=\nu((0,s])\). Assume
\begin{equation}
             U(s)\sim s^\alpha L(1/s),\qquad s\downarrow0,
             \quad \alpha\ge0,
 \label{eq:rv}
\end{equation}
where \(L\) is slowly varying and positive at infinity.

\begin{theorem}[Balanced time and spatial dilation]
\label{thm:balanced}
If \(R\to\infty\) and \(b\to b_0\in(0,\infty)\), then
\[
 V(R,\tau)\sim U(R^{-2})\,\mathcal M_{b_0}(\alpha),\qquad
 \mathcal M_b(\alpha)=
 \begin{cases}
 \displaystyle\alpha\int_0^\infty F_b(y)y^{\alpha-1}\dd y,
             &\alpha>0,\\[2pt]
 F_b(0),&\alpha=0.
 \end{cases}
\]
The constants are finite and strictly positive. There is no
saturation of the threshold index at \(1/2\) in this experiment.
\end{theorem}

\begin{proof}
Set \(X=R^2\) and \(m_X(A)=\nu(A/X)/U(1/X)\).
The cumulative ratios give vague convergence on \([0,\infty)\)
to \(\alpha y^{\alpha-1}\dd y\) for \(\alpha>0\) and
\(\delta_0\) for \(\alpha=0\).
By Potter bounds \cite[Thm.~1.5.6]{BGT},
\(m_X([0,y])\le C_\delta y^{\alpha+\delta}\) for
\(1\le y\le s_* X\), with \(s_*>0\) sufficiently small.
The remaining integral \(y>s_*X\) against a rapidly decaying
kernel is negligible by the original polynomially weighted
finite tail and slow variation of \(U(1/X)\).
Dyadic tail estimates and local uniform kernel convergence
now give the formula. At \(\alpha=0\), an atom-free continuum
with slowly vanishing cumulative also has a zero-scale
tangent atom; its absolute variance still tends to zero.
\end{proof}

\subsection{Short and long time gates}

\begin{theorem}[Anisotropic asymptotics]
\label{thm:anisotropic}
Keep \eqref{eq:rv} and take \(R\to\infty\).
\begin{enumerate}
\item If \(b\to0\), \(\tau\) is bounded away from zero, and
\(0\le\alpha<1/2\), the formula of Theorem~\ref{thm:balanced}
holds with \(b_0=0\) and \(F_0\) given by \eqref{eq:scale}
with the last factor replaced by \(1\).
\item If \(b\to0\), \(\tau\to\infty\), and \(\alpha>1/2\), then
\begin{equation}
 V(R,\tau)\sim R^{-1}\tau^{1-2\alpha}L(\tau^2)
       \frac{\alpha W_4}{2}\int_0^\infty
                     y^{\alpha-3/2}|\widehat g(\sqrt y)|^2\dd y .
 \label{eq:shortsuper}
\end{equation}
\item If \(b\to0\), \(\tau\to\infty\), and \(\alpha=1/2\), then
\begin{equation}
       V(R,\tau)\sim\frac{W_4}{4R}
                     \int_{\tau^2}^{R^2} L(t)\,\frac{\dd t}{t}.
 \label{eq:shortcritical}
\end{equation}
For \(L(t)\to c>0\), this is
\(cW_4(2R)^{-1}\log(R/\tau)\).
\item If \(b\to\infty\) (so \(\tau\to\infty\)), let \(P_m\) be the first nonzero
homogeneous Fourier jet of \(w\) at zero as above. Then
\begin{equation}
 V(R,\tau)\sim b^{-6-2m}U(\tau^{-2})\,\mathcal M_*(\alpha),
 \label{eq:long}
\end{equation}
where \(\mathcal M_*\) is defined as in Theorem~\ref{thm:balanced}
using the kernel
\[
 F_*(y)=\frac12\int\frac{\dd^3\boldsymbol l}{(2\pi)^3}
 \frac{|\boldsymbol l|^4|P_m(\boldsymbol l)|^2}
           {\sqrt{|\boldsymbol l|^2+y}}
          |\widehat g(\sqrt{|\boldsymbol l|^2+y})|^2 .
\]
Even an actual null atom is suppressed by the factor
\(b^{-6-2m}\) in this regime.
\end{enumerate}
\end{theorem}

\begin{proof}
For the first case the kernels converge locally to \(F_0\)
and are bounded by \(C(1+y)^{-1/2}\) uniformly in \(b\).
On \(1\le y\le s_*R^2\), choose \(\delta\) with
\(\alpha+\delta<1/2\) in the dyadic estimate in the previous
proof. On the original tail \(s>s_*\) the normalized temporal
gate with \(\tau\) bounded below gives an integrable bound
\(C_NR^{-1}s^{-1/2}(1+s)^{-N}\) in \eqref{eq:scale},
hence an \(O(R^{-1})=o(U(R^{-2}))\) contribution. This also
shows finiteness of the limiting integral with \(F_0\).

For \eqref{eq:shortsuper}, set \(y=\tau^2s\) and write
\(F_b(y/b^2)=bK_b(y)\), where
\[
 K_b(y)=\frac12\int\frac{\dd^3\mathbf k}{(2\pi)^3}
 |\mathbf k|^4|\widehat w(\mathbf k)|^2
 \frac{|\widehat g(\sqrt{b^2|\mathbf k|^2+y})|^2}
            {\sqrt{b^2|\mathbf k|^2+y}}.
\]
This tends on \(y>0\) to \(W_4|\widehat g(\sqrt y)|^2/(2\sqrt y)\)
and is uniformly bounded by \(C_N y^{-1/2}(1+y)^{-N}\).
Near zero, Potter gives
\(\nu([0,y/\tau^2])/U(\tau^{-2})\le C_\delta
 y^{\alpha-\delta}\) with \(\alpha-\delta>1/2\), which
controls the inverse square root. At infinity use the rapid
bound as before. The limiting integral at scale \(\tau^2\)
and the extra factor \(b=\tau/R\) give \eqref{eq:shortsuper}.

For the critical case fix \(K>1\). On
\(K/R^2\le s\le 1/(K\tau^2)\) the ratio of the kernel to
\(W_4/(2R\sqrt s)\) is uniformly \(1+o_K(1)\), with the
error tending to zero as \(K\to\infty\). To see this truncate
to a compact \(\mathbf k\)-set using the Schwartz spatial
tail, take \(R^2s\ge K\) in the denominator and use
\(\tau^2s\le1/K\) and \(b\to0\) in the temporal gate.
The lower and upper discarded ranges contribute,
respectively, \(O_K(R^{-1}L(R^2))\) and
\(O_K(R^{-1}L(\tau^2))\).
The first estimate uses the bounded \(F_b\) and the cumulative
measure at \(K/R^2\); the second uses \(bK_b(y)\) for
\(y\ge1/K\) and rapid decay. Integrating \(s^{-1/2}\)
against the cumulative \(U(s)\) between the two cutoffs gives
\(\frac12\int L(1/s)\dd s/s\), up to endpoint terms of the
same two orders before the common \(R^{-1}\).
The asymptotic assumption on \(U\) is uniform once both
endpoints are small. Finally
\(\int_{\tau^2}^{R^2}L(t)\dd t/t\)
dominates \(L(\tau^2)+L(R^2)\): an arbitrary fixed logarithmic
length at each end is eventually present and yields that
length times its endpoint value by uniform convergence of
slow variation. Thus the discarded terms and fixed-width
strips are negligible for each \(K\); letting \(K\to\infty\)
gives \eqref{eq:shortcritical}.

For the last case substitute \(\mathbf k=\boldsymbol l/b\)
in \eqref{eq:scale}. The spatial powers and denominator
give \(b^{-6}\); the remaining profile
\(b^m\widehat w(\boldsymbol l/b)\) tends to \(P_m(\boldsymbol l)\)
and is bounded globally by \(C|\boldsymbol l|^m\), by Taylor
expansion near zero and boundedness of \(\widehat w\) away
from zero. This gives the extra \(b^{-2m}\) and normalized
kernels uniformly rapid in \(\tau^2s\). The proof of the
balanced limit at scale \(\tau^2\) applies.
\end{proof}

If \(\tau\) is fixed and the nonzero-mass part satisfies
\(\int_{(0,1)}s^{-1/2}\nu(\dd s)<\infty\), then \(R\)
times its variance converges to
\[
 \frac{W_4}{2}\int_{(0,\infty)}
       s^{-1/2}|\widehat g(\tau\sqrt s)|^2\nu(\dd s),
\]
by dominated convergence in the shell formula. This sees
the full gated mass rather than only the low-end cumulative
measure, and its coefficient can vanish at temporal notches.
In particular an atom of weight \(A>0\) at \(s=m^2>0\), with a
zero of order \(h\) of \(\widehat g\) at \(\tau m\) (order zero
means nonzero), contributes precisely
\[
 \frac{A W_{4+4h}}{2m}
  \left|\frac{\widehat g^{(h)}(\tau m)}{h!}\right|^2
  \left(\frac{\tau}{2m}\right)^{2h} R^{-1-4h}(1+o(1)).
\]
Indeed the departure of the gate's argument from \(\tau m\)
is \(\tau|\mathbf k|^2/(2mR^2)+O(R^{-4})\) at fixed
\(\mathbf k\), and a global bound by the leading power in
\(|\mathbf k|/R\) gives domination by the Schwartz window.

For the weighted spin measures of Proposition~\ref{thm:spin},
the analogous formulas replace the spatial factor \(q^4\)
by \(1\). Removing the exact prefactor \(R^4\) gives
\eqref{eq:scale} with \(|\mathbf k|^4\) omitted; the long-time
exponent on this normalized variance is then \(2+2m\),
not \(6+2m\), from the same change of variables. If an
unweighted spin cumulative has positive index \(\beta\),
weighting it by \(s^2\) changes the index to \(\beta+2\), with
cumulative factor \(\beta/(\beta+2)\) by integration by parts.
The index-zero case requires its own slow-variation data.

\section{A Tauberian converse using delay ratios}
\label{sec:converse}

A measured log-slope alone does not force a regularly varying
threshold. For example, for \(\alpha>0\) a sufficiently small
positive log-periodic modulation of \(s^{\alpha-1}\dd s\)
near zero gives a non-regularly-varying cumulative while
preserving the log-slope of \eqref{eq:scale}. In general
inverse regular variation through a single transform depends
on its Mellin cancellation properties \cite{JMRS}. The
positive delay family gives a converse without a zero-free
assumption on one Mellin kernel.

For example, if a single chosen kernel happened to satisfy
\(\int_0^\infty F_1(y)y^{\alpha+i\beta-1}\dd y=0\) at
\(\alpha>0\), \(\beta\ne0\), consider the positive measure
\(s^{\alpha-1}[1+\delta\cos(\beta\log s)]\dd s\) on \((0,1)\),
\(0<\delta<1\), in fixed units. Its cumulative divided by
\(s^\alpha\) has a nonconstant log-periodic part, since the
integral of the oscillatory density contributes
\(\operatorname{Re}[s^{i\beta}/(\alpha+i\beta)]\).
However its reference variance under \(y=Xs\) has leading
term \(X^{-\alpha}\int F_1(y)y^{\alpha-1}\dd y\): the leading
oscillation vanishes by the assumed Mellin zero and the
finite upper cutoff contributes only a rapid error. This
conditional obstruction explains why the proof below brings
in a separating family; we do not assert that a given
physical choice of \(w,g\) has such a zero.

Fix \(w,g\) and use balanced cutoff \(R\) for both scales.
At scale \(X=R^2\), let \(V_\theta(X)\) denote the energy
variance with temporal shape \(g^{[\theta]}\) dilated by
\(R\), i.e. the laboratory delay is \(\theta R\).
The reference is \(A_X=V_0(X)\) (or its atom-subtracted
version, if all tested data have also had that atom term
removed).

\begin{theorem}[Delay-family Tauberian equivalence]
\label{thm:converse}
Suppose the reference \(A_X>0\) is regularly varying at
infinity with index \(-\alpha\) and the ratios
\(V_\theta(X)/A_X\) have limits on a dense set of nonnegative
\(\theta\). Then \(\alpha\ge0\) and the corresponding mass
cumulative satisfies
\[
             U(s)\sim
            \frac{A_{1/s}}{\mathcal M_1(\alpha)},\qquad s\downarrow0.
\]
Conversely \eqref{eq:rv} gives regular variation of the
reference and all of these ratio limits.
\end{theorem}

\begin{proof}
Boundedness of the balanced variance implies
\(\alpha\ge0\). Put \(\mu_X(B)=\nu(B/X)/A_X\).
Strict positivity of the reference kernel on \([0,1]\)
gives \(\nu([0,t^{-1}])\le C A_t\) for all large \(t\).
There is therefore uniform mass control on compact sets,
and for \(1\le y\le X/X_0\) Potter applied to
\(A_{X/y}/A_X\) gives
\(\mu_X([0,y])\le C_\delta y^{\alpha+\delta}\).
On the remaining tail the original polynomially
weighted spectral bound suffices against the rapid
test, after dividing by \(A_X\).
The delay kernels have the same uniform tail bound because
\(|\widehat{g^{[\theta]}}|\le|\widehat g|\).
Any vague subsequential limit is thus a locally finite
positive measure of polynomial growth with
\(\int F_1\,\dd\mu=1\) and the same family of limiting
integrals. Its delay transform is continuous, so
uniqueness for the dense set extends to all delays.

Apply Theorem~\ref{thm:inverse} to these tangent measures
at fixed dimensionless profiles: its proof only requires
local finiteness and polynomially weighted integrability,
which the growth bound gives. Hence all subsequential
limits coincide. The identity
\[
 \mu_{cX}(B)=(A_X/A_{cX})\,\mu_X(B/c),\qquad c>0,
\]
then forces \(\mu([0,y])=\mu([0,1])y^\alpha\).
For \(\alpha=0\) this is a delta at the origin; otherwise
it is the homogeneous power-law measure. Testing with
\(F_1\) fixes \(\mu([0,1])=1/\mathcal M_1(\alpha)\).
Its boundary at \(1\) has zero mass, proving the stated
cumulative equivalence. The converse is Theorem~\ref{thm:balanced}
applied to each compact double pulse separately.
\end{proof}

The fixed-profile inverse and the scaling equivalence answer
different questions. Exact lag data reconstruct the positive
mass measure, but Section~\ref{sec:noise} prevents uniform
fixed-window mass discrimination. The scaling asymptotics
separate a null atom from a slowly vanishing continuum only
when its absolute rather than its normalized limiting
amplitude is used. The results do not identify a stress spectrum
with a scattering pole measure or determine the nonlinear
gravity fluctuations of a total charge.

\begin{thebibliography}{99}
\small
\bibitem{P45}
A massless stress pole from the variance of a spatially scaled vacuum energy,
companion manuscript 45,
\href{https://github.com/g3173028019-beep/ai-automated-quantum-gravity-research/tree/main/papers/45-vacuum-stress-massless-pole}{repository version}.
\bibitem{CFL}
A.~Cappelli, D.~Friedan, and J.~I.~Latorre,
C-theorem and spectral representation,
\emph{Nucl. Phys. B} \textbf{352} (1991), 616--670,
\href{https://doi.org/10.1016/0550-3213(91)90102-4}{doi:10.1016/0550-3213(91)90102-4}.
\bibitem{Karateev}
D.~Karateev, Two-point functions and bootstrap applications in quantum field theories,
\emph{JHEP} \textbf{02} (2022), 186,
\href{https://arxiv.org/abs/2012.08538}{arXiv:2012.08538}.
\bibitem{Karamata}
J.~Karamata, Neuer Beweis und Verallgemeinerung der Tauberschen S\"atze, welche die Laplacesche und Stieltjessche Transformation betreffen,
\emph{J. Reine Angew. Math.} \textbf{164} (1931), 27--39,
\href{https://geodesic.mathdoc.fr/item/JRAM_1931__164_149745/}{archive edition}.
\bibitem{BGT}
N.~H.~Bingham, C.~M.~Goldie, and J.~L.~Teugels,
\emph{Regular Variation}, Encyclopedia of Mathematics and its Applications, vol.~27, Cambridge University Press (1987),
\href{https://doi.org/10.1017/CBO9780511721434}{doi:10.1017/CBO9780511721434}.
\bibitem{JMRS}
M.~Jacobsen, T.~Mikosch, J.~Rosi\'nski, and G.~Samorodnitsky,
Inverse problems for regular variation of linear filters, a cancellation property for \(\sigma\)-finite measures, and identification of stable laws,
\emph{Ann. Appl. Probab.} \textbf{19} (2009), 210--242,
\href{https://arxiv.org/abs/0712.0576}{arXiv:0712.0576}.
\end{thebibliography}
\end{document}
